\documentclass[aspectratio=169,t]{beamer}

% ---- Theme ----
\usetheme{Madrid}
\usecolortheme{whale}
\setbeamertemplate{navigation symbols}{}
\setbeamertemplate{footline}[frame number]
\setbeamerfont{footnote}{size=\tiny}

% ---- Packages ----
\usepackage{booktabs}
\usepackage{colortbl}
\usepackage{multirow}
\usepackage{tikz}
\usetikzlibrary{arrows.meta,backgrounds,positioning,fit,calc}
\usepackage{amssymb}
\usepackage{enumitem}
\usepackage{xcolor}

% ---- Colors (match paper) ----
\definecolor{c1blue}{HTML}{4C72B0}
\definecolor{c2orange}{HTML}{DD8452}
\definecolor{neutralgray}{HTML}{666666}
\definecolor{evalgreen}{HTML}{55A868}
\definecolor{bgoffline}{HTML}{EEF2F8}
\definecolor{bgonline}{HTML}{FEF6EF}
\definecolor{highlight}{HTML}{E74C3C}

% ---- Metadata ----
\title[LectureBench]{LectureBench:\\Long-Lecture Video RAG Benchmark}
\subtitle{Temporally Grounded QA over Academic Lectures}
\author{Nasim Faridnia}
\institute{Department of Computer Science\\University of Texas at San Antonio}
\date{2026}

% ---- Helper: table row highlight ----
\newcommand{\rowhl}{\rowcolor{yellow!25}}

%=============================================================
\begin{document}
%=============================================================

%------------------------------------------------------------
\begin{frame}
\titlepage
\note{Introduce the paper title and context. This is a benchmark paper for
BMVC 2026 that proposes LectureBench — a temporally grounded QA dataset and
RAG evaluation framework over long academic lecture videos.}
\end{frame}

%------------------------------------------------------------
\begin{frame}{The Problem: Long Lectures Are Not Served by Existing Benchmarks}

\begin{columns}[T]
\begin{column}{0.48\textwidth}
  \textbf{Short-clip benchmarks dominate VideoQA}
  \begin{itemize}
    \item ActivityNet-QA, NExT-QA: avg $<$3 min per clip
    \item EgoSchema: 3-minute egocentric clips
    \item Optimized for moment classification, not evidence retrieval
  \end{itemize}

  \medskip
  \textbf{Lecture videos are fundamentally different}
  \begin{itemize}
    \item \textbf{47--89 min} per lecture
    \item Answers may span multiple evidence segments \textbf{60+ minutes apart}
    \item Critical information lives on \textbf{slides and whiteboards} — absent from spoken transcript
  \end{itemize}
\end{column}
\begin{column}{0.48\textwidth}
  \textbf{Current lecture benchmarks fall short}
  \begin{itemize}
    \item EduVidQA: fixed 4-min context window — retrieval problem sidestepped entirely
    \item No existing benchmark evaluates \textbf{temporal grounding} in lecture RAG
    \item No benchmark requires multi-hop evidence spanning non-adjacent segments
  \end{itemize}

  \bigskip
  \begin{alertblock}{The Gap}
    No benchmark evaluates whether a RAG system retrieves the \emph{right moment}
    from a long lecture — not merely a semantically similar one.
  \end{alertblock}
\end{column}
\end{columns}

\note{Emphasize: the core gap is temporal precision + multi-hop + visual grounding
combined on lecture-length content. Existing work covers each property in isolation
or on short clips.}
\end{frame}

%------------------------------------------------------------
\begin{frame}{Comparison with Related Benchmarks}

\begin{center}
\small
\begin{tabular}{lcccccc}
\toprule
Dataset & Domain & \#Videos & \#QA & Avg Dur. & Multi-hop & Temporal \& Visual \\
\midrule
NExT-QA           & General     & 5,440 & 52,044 & 44s     & $\times$ & $\times$ \\
EgoSchema         & Egocentric  & —     &  5,031 & 3 min   & $\times$ & $\times$ \\
NA-VQA            & Movies      &    88 &  4,442 & 80--200 min & implicit & $\times$ \\
LongVidSearch     & Mixed       &   447 &  3,000 & 26 min  & \checkmark & $\times$ \\
EduVidQA          & Lectures    &   296 &  5,252 & 68 min  & $\times$ & \checkmark \\
\midrule
\rowhl\textbf{LectureBench (ours)} & \textbf{Lectures} & \textbf{60} & \textbf{810}
  & \textbf{63 min} & \checkmark & \checkmark \\
\bottomrule
\end{tabular}
\end{center}

\medskip
\begin{columns}[T]
\begin{column}{0.3\textwidth}
  \textbf{Gap 1:} No benchmark is multi-hop \emph{and} lecture-domain
\end{column}
\begin{column}{0.3\textwidth}
  \textbf{Gap 2:} No benchmark anchors answers to timestamps \emph{and} has a visual split
\end{column}
\begin{column}{0.3\textwidth}
  \textbf{Gap 3:} EduVidQA sidesteps retrieval with fixed 4-min windows
\end{column}
\end{columns}

\note{[TABLE] This reproduces Table 1 from the paper (benchmark\_comparison).
You can use the paper's Table 1 directly as a high-resolution screenshot or
re-typeset it on the slide.

Key talking point: LectureBench is the ONLY benchmark that combines multi-hop,
temporal span evaluation, visual-dependent questions, and lecture-length videos.}
\end{frame}

%------------------------------------------------------------
\begin{frame}{Contributions}

\begin{enumerate}
  \setlength{\itemsep}{10pt}
  \item \textbf{Corpus:} 60 CC-licensed academic lecture videos (62.6 hours) across
        15 courses from MIT OCW, NYU Deep Learning, and Stanford SEE.

  \item \textbf{Benchmark:} 810 LLM-reviewed QA pairs with precise temporal ground-truth
        spans, typed by \emph{multi-hop-visual}, \emph{visual}, \emph{multi-hop},
        \emph{text}, and \emph{unanswerable} categories.

  \item \textbf{Open-source RAG pipeline} with two retrieval configurations and a
        comprehensive evaluation suite measuring temporal IoU, Hit Rate@$k$,
        LLM-judge score, and citation accuracy.

  \item \textbf{Empirical finding:} visual frame captions yield \textbf{+65\% tIoU}
        on visual-dependent questions but are \textbf{essentially unchanged} on
        text-only questions, precisely isolating the multimodal contribution.
\end{enumerate}

\note{Contribution 4 is the central empirical result — frame captions help exactly
where they should (visual evidence) and do not hurt where they are irrelevant
(transcript-only questions). This is the paper's punchline.}
\end{frame}

%------------------------------------------------------------
\begin{frame}{LectureBench at a Glance}

\begin{columns}[T]
\begin{column}{0.44\textwidth}
  \small
  \begin{tabular}{lc}
  \toprule
  Statistic & Value \\
  \midrule
  Lectures                      & 60 \\
  Sources                       & MIT OCW, NYU DL, Stanford SEE \\
  Total duration                & 62.6 hours \\
  Duration range                & 47--89 min (avg 63 min) \\
  Transcript chunks (45s/35s)   & 7,382 \\
  Frames captioned (every 15s)  & 17,180 \\
  QA pairs (accepted)           & 810 (avg 13.5/lecture) \\
  Visual-dependent pairs        & 56.2\% (455/810) \\
  \bottomrule
  \end{tabular}

  \medskip\footnotesize
  Dataset released under CC BY-NC-SA 4.0.\\
  All code and data: \texttt{github.com/NassimF/VQA-Benchmark}
\end{column}
\begin{column}{0.52\textwidth}
  \textbf{Five question types:}
  \medskip
  \small
  \begin{tabular}{lcl}
  \toprule
  Type & Count & Role \\
  \midrule
  Multi-hop visual & 320 & Primary \\
  Single-hop visual & 135 & Primary \\
  Multi-hop textual & 140 & Control \\
  Single-hop textual & 103 & Control \\
  Unanswerable & 112 & Required \\
  \midrule
  \textbf{Total} & \textbf{810} & \\
  \bottomrule
  \end{tabular}

  \medskip
  \textbf{Visual-dependent} (multi-hop-visual + visual) = 455 pairs (56.2\%)\\
  \textbf{Text-only control} = 243 pairs (30\%)
\end{column}
\end{columns}

\note{[TABLE] This reproduces Table 2 (corpus stats) from Section 4.1 of the paper.
You can paste Table 2 on the left column as a figure for higher fidelity.

[FIGURE] Optionally add Figure 6 (the QA distribution bar chart, fig\_qa\_dist) on
the right side instead of the manual table — it gives a cleaner visual breakdown.}
\end{frame}

%------------------------------------------------------------
\begin{frame}{End-to-End Pipeline (Figure 2)}

\tikzset{
  box/.style={rectangle, rounded corners=3pt, draw=#1, fill=#1!10,
              text width=18mm, align=center, font=\small, inner sep=4pt,
              minimum height=9mm},
  cfgbox/.style={rectangle, rounded corners=3pt, draw=#1, fill=#1!15,
              text width=22mm, align=center, font=\small, inner sep=4pt,
              minimum height=9mm},
  evalbox/.style={rectangle, rounded corners=3pt, draw=evalgreen, fill=evalgreen!10,
              text width=24mm, align=center, font=\small, inner sep=4pt,
              minimum height=9mm},
  arr/.style={-{Stealth[length=5pt]}, semithick, #1},
  seclabel/.style={font=\small\bfseries, text=#1},
}

\begin{center}
\resizebox{0.96\textwidth}{!}{%
\begin{tikzpicture}[node distance=7mm and 9mm]
  \node[box=neutralgray] (video)  {Lecture Video};
  \node[box=neutralgray, right=of video]  (vtt)   {VTT Transcript};
  \node[box=neutralgray, right=of vtt]    (chunk) {Chunker\\{\small 45s/35s stride}};
  \node[cfgbox=c1blue,   right=of chunk]  (db1)
        {\textbf{\textcolor{c1blue}{Config~1}}\\{\small transcript only}\\[-1pt]
         {\scriptsize\textcolor{c1blue!60!black}{\textit{ChromaDB}}}};
  \node[box=neutralgray, right=22mm of db1] (retriever)
        {Retriever\\{\small top-$k$, cosine}};
  \node[box=neutralgray, right=of retriever] (gen)
        {Generator\\{\small GPT-4o-mini}};
  \node[box=neutralgray, right=of gen]       (answer)
        {Answer +\\Citations};

  \node[cfgbox=c2orange, below=16mm of chunk] (cap)
        {\textbf{\textcolor{c2orange}{Frame Caps}}\\{\small Qwen2-VL-7B}};
  \node[cfgbox=c2orange] (db2) at (db1 |- cap)
        {\textbf{\textcolor{c2orange}{Config~2}}\\{\small text + captions}\\[-1pt]
         {\scriptsize\textcolor{c2orange!60!black}{\textit{ChromaDB}}}};
  \node[evalbox] (gt)      at (retriever |- cap) {GT Spans\\{\small (benchmark)}};
  \node[evalbox] (eval)    at (gen      |- cap)  {Evaluator};
  \node[evalbox] (metrics) at (answer   |- cap)  {tIoU $\cdot$ HR@$k$\\LLM-judge $\cdot$ Cit.~Acc.};

  \draw[arr=neutralgray] (video) -- (vtt);
  \draw[arr=neutralgray] (vtt)   -- (chunk);
  \draw[arr=c1blue] (chunk) -- (db1);
  \draw[arr=c1blue] (db1)   -- (retriever);
  \draw[arr=c2orange] (video.south) |- (cap.west);
  \draw[arr=c2orange] (cap) -- (db2);
  \coordinate (c2exit) at ($(db1.east)!0.6!(retriever.west)$);
  \draw[arr=c2orange] (db2.east) -- (c2exit |- db2) |- (retriever.west);
  \draw[arr=neutralgray] (retriever) -- (gen);
  \draw[arr=neutralgray] (gen)       -- (answer);
  \draw[arr=evalgreen] (answer.south) -- (eval.north east);
  \draw[arr=evalgreen] (gt)  -- (eval);
  \draw[arr=evalgreen] (eval) -- (metrics);

  \begin{scope}[on background layer]
    \node[fill=bgoffline, rounded corners=5pt, inner sep=6pt,
          fit=(video)(vtt)(chunk)(db1)(cap)(db2)] (offbox) {};
    \node[fill=bgonline,  rounded corners=5pt, inner sep=6pt,
          fit=(retriever)(gen)(answer)(gt)(eval)(metrics)] (onbox) {};
  \end{scope}
  \node[seclabel=neutralgray, above=2mm of offbox.north west, anchor=south west]
        {Offline Indexing};
  \node[seclabel=neutralgray, above=2mm of onbox.north west, anchor=south west]
        {Online Query \& Evaluation};
\end{tikzpicture}%
}
\end{center}

\note{This is Figure 2 from the paper, reproduced directly as TikZ.

Key points to narrate:
- LEFT (blue background): offline preprocessing — same for both configs except
  Config 2 appends Qwen2-VL-7B frame captions before indexing.
- RIGHT (orange background): both configs share the SAME retriever, generator,
  and evaluator — so any difference in results is purely due to chunk content.
- The controlled comparison is the paper's main methodological contribution.}
\end{frame}

%------------------------------------------------------------
\begin{frame}{Pipeline: Chunking \& Frame Captioning}

\begin{columns}[T]
\begin{column}{0.48\textwidth}
  \textbf{Temporal Chunking}
  \begin{itemize}
    \item \textbf{Window:} 45 seconds / \textbf{Stride:} 35 seconds (10s overlap)
    \item 45s $\approx$ one complete concept explanation per lecture
    \item Overlap ensures evidence near chunk boundaries appears in both neighbours
    \item Result: \textbf{7,382 chunks} across 60 lectures
  \end{itemize}

  \medskip
  \textbf{Why not shot-boundary chunking?}\\
  Lecture recordings have no shot cuts — fixed-width windows are the only
  principled alternative and keep the pipeline interpretable.
\end{column}
\begin{column}{0.48\textwidth}
  \textbf{Frame Captioning (Config 2 only)}
  \begin{itemize}
    \item One frame extracted every \textbf{15s} via \texttt{ffmpeg}
    \item Captioned by \textbf{Qwen2-VL-7B-Instruct} (bfloat16, single A100)
    \item Prompt: describe all visible text, equations, diagrams in detail
    \item Caption appended as \texttt{[frame caption: \ldots]} to chunk text
    \item \textbf{17,180 frames} total; $\approx$3 frames per chunk
    \item Cost: $\approx$2--8 GPU-hours for all 60 lectures
  \end{itemize}

  \medskip
  \begin{exampleblock}{Key Design Choice}
    Appending captions as text keeps the \textbf{embedding model fixed}
    across both configs — differences in retrieval quality are due to chunk
    content alone, not embedding architecture.
  \end{exampleblock}
\end{column}
\end{columns}

\note{The 15s frame interval guarantees at least one frame per distinct slide
state within any 45s chunk (slides typically displayed 30--120s).

Emphasize: fixing all-MiniLM-L6-v2 across both configs is a deliberate
controlled-experiment decision. The only variable is what text is embedded.}
\end{frame}

%------------------------------------------------------------
\begin{frame}{Pipeline: Retrieval \& Answer Generation}

\begin{columns}[T]
\begin{column}{0.48\textwidth}
  \textbf{Retrieval}
  \begin{itemize}
    \item Embedding: \textbf{all-MiniLM-L6-v2} (384-dim, Sentence-BERT)
    \item Vector store: \textbf{ChromaDB} — two persistent collections:
      \begin{itemize}
        \item \textcolor{c1blue}{\texttt{transcript\_only}}
        \item \textcolor{c2orange}{\texttt{transcript\_plus\_frames}}
      \end{itemize}
    \item Query-time: top-$k$ cosine similarity, $k \in \{1, 3, 5\}$
    \item Same model + store for both configs — retrieval difference is \emph{only} chunk content
  \end{itemize}
\end{column}
\begin{column}{0.48\textwidth}
  \textbf{Answer Generation}
  \begin{itemize}
    \item Generator: \textbf{GPT-4o-mini}, temperature 0.1
    \item Input: retrieved chunks as numbered excerpts
    \item Output: structured JSON — \texttt{answer} + \texttt{citations} list
    \item Predicted span = \textbf{union} of all cited chunk time windows
    \item Citations rendered as \texttt{[video\_id @ mm:ss to mm:ss]}
  \end{itemize}

  \medskip
  \begin{alertblock}{Why a weak generator?}
    A lightweight generator isolates retrieval failures that a stronger model might
    mask by drawing on world knowledge outside the retrieved chunks.
  \end{alertblock}
\end{column}
\end{columns}

\note{The choice of GPT-4o-mini as generator (rather than GPT-4o or Claude) is
intentional and should be emphasized: the variable under study is RETRIEVAL CONFIG,
not generator quality. A stronger generator would compress the gap between configs
by hallucinating correct answers from parametric memory.}
\end{frame}

%------------------------------------------------------------
\begin{frame}{Benchmark Construction: QA Generation}

\begin{columns}[T]
\begin{column}{0.48\textwidth}
  \textbf{Generation}
  \begin{itemize}
    \item \textbf{Claude Sonnet 4.6} generates \textbf{15 draft QA pairs} per lecture
          from augmented chunks (transcript + frame captions)
    \item 5 types targeted per lecture (7 multi-hop-visual, 2 visual, 3 multi-hop,
          2 text, 1 unanswerable)
    \item Multi-hop pairs must cite $\geq$2 \textbf{non-adjacent} spans
          ($\geq$70s gap required)
    \item Total cost: $\approx$\$8 for all 60 lectures
  \end{itemize}

  \medskip
  \begin{exampleblock}{Generator $\neq$ Evaluator}
    QA generation uses Claude Sonnet 4.6; answer generation at eval time uses
    GPT-4o-mini. Using the same model for both would be circular.
  \end{exampleblock}
\end{column}
\begin{column}{0.48\textwidth}
  \textbf{LLM Review (two-pass filter)}
  \begin{itemize}
    \item \textbf{Pass 1 — Claude Sonnet 4.6:} three-criterion binary rubric:
      \begin{itemize}
        \item C1: answer correctness vs.\ chunk text
        \item C2: span plausibility ($\geq$70s gap)
        \item C3: question type accuracy
      \end{itemize}
    \item \textbf{Pass 2 — Cross-family check:} pairs passing C1 re-evaluated by
          GPT-4o on C1 alone; disagreement triggers rejection
    \item Atomic fact decomposition (FActScore) before C1 verification
    \item Acceptance rate: $\approx$90\% (810 accepted / 900 target)
  \end{itemize}
\end{column}
\end{columns}

\note{The cross-family C1 check (Claude + GPT-4o) guards against knowledge-conflict
failures where a judge accepts a claim based on its own parametric knowledge rather
than the chunk evidence. This is the key quality-control innovation.}
\end{frame}

%------------------------------------------------------------
\begin{frame}{Example QA Pair: Multi-hop Visual}

\begin{center}
\small
\setlength{\fboxsep}{8pt}
\fbox{
\begin{minipage}{0.90\textwidth}
\textbf{Question (multi-hop-visual):}
What does the blackboard show as the theme of today's lecture, and how does the
instructor connect this theme to the interval scheduling problems he discusses?

\medskip
\textbf{Ground-truth spans:}
\begin{itemize}[noitemsep,topsep=2pt,leftmargin=*]
  \item \textbf{Hop 1} [10:20--11:15] \emph{(visual)} ---
        \texttt{[frame caption: Blackboard shows ``Theme of today's lecture ---
        Very similar problems can have very different complexity.'']}
  \item \textbf{Hop 2} [79:35--80:40] \emph{(transcript)} ---
        ``\ldots interval scheduling is greedy $O(n)$; weighted interval scheduling
        requires $O(n^2)$ DP; non-identical machines is NP-complete \ldots''
\end{itemize}

\textbf{Span gap:} $4775 - 675 = 4100$\,s (68 min) $\gg$ 70\,s required.
\quad \textbf{Source:} \texttt{mit\_6046\_lec01\_q017}
\end{minipage}
}
\end{center}

\note{[FIGURE] This is a simplified reproduction of Figure 3 from the paper
(fig:example\_qa). For the advisor presentation, consider adding the timeline
diagram from Figure 3 showing Hop 1 at 10:20 and Hop 2 at 79:35 with the
``68 min gap'' annotation — it visually makes the retrieval challenge clear.

Key point: Hop 1's answer exists ONLY in the frame caption (blackboard text),
not in the spoken transcript. Config 1 would retrieve the right time window
but fail to answer because the evidence is invisible without frame captions.}
\end{frame}

%------------------------------------------------------------
\begin{frame}{Question Type Distribution}

\begin{columns}[T]
\begin{column}{0.50\textwidth}
  \small
  \begin{tabular}{lcc}
  \toprule
  Type & Count & \% \\
  \midrule
  Multi-hop visual   & 320 & 39.5\% \\
  Single-hop visual  & 135 & 16.7\% \\
  \midrule
  \textbf{Visual-dependent} & \textbf{455} & \textbf{56.2\%} \\
  \midrule
  Multi-hop textual  & 140 & 17.3\% \\
  Single-hop textual & 103 & 12.7\% \\
  \textbf{Text-only} & \textbf{243} & \textbf{30.0\%} \\
  \midrule
  Unanswerable & 112 & 13.8\% \\
  \midrule
  \textbf{Total} & \textbf{810} & \textbf{100\%} \\
  \bottomrule
  \end{tabular}

  \medskip
  \begin{itemize}
    \item Multi-hop questions: \textbf{460/810 (56.8\%)}
    \item Target was $\geq$60\% visual; achieved \textbf{56.2\%}
    \item Gap reflects higher difficulty of $\geq$70s span requirement
  \end{itemize}
\end{column}
\begin{column}{0.46\textwidth}
  \begin{center}
  \vspace{4pt}
  \fbox{\rule{0pt}{4.8cm}\rule{5.5cm}{0pt}}
  \end{center}
  \vspace{-4pt}
  \footnotesize\centering
  \textit{[Insert Figure 6 — fig\_qa\_dist bar chart here]}
\end{column}
\end{columns}

\note{[FIGURE] Place Figure 6 (images/fig\_qa\_dist.pdf) in the right column
placeholder. The bar chart groups pairs into visual-dependent vs control categories
and gives an immediate visual read of the 56\% / 44\% split.

Design rationale to mention: the text-only control group (43.8\%) is essential for
isolating Config 2's improvement — without it, we couldn't tell whether gains came
from frame captions or from generally better embeddings.}
\end{frame}

%------------------------------------------------------------
\begin{frame}{Evaluation Metrics}

\begin{columns}[T]
\begin{column}{0.48\textwidth}
  \textbf{Temporal Grounding}
  \begin{itemize}
    \item \textbf{Mean Temporal IoU} — overlap between predicted span (union of cited chunks)
          and ground-truth span
    \item \textbf{IoU@0.3} — fraction of pairs with tIoU $\geq$ 0.3
    \item \textbf{IoU@0.5} — fraction with tIoU $\geq$ 0.5 (stricter)
  \end{itemize}

  \medskip
  \textbf{Retrieval Coverage}
  \begin{itemize}
    \item \textbf{Hit Rate@$k$} ($k \in \{1,3,5\}$) — fraction where ground-truth span
          is covered by at least one top-$k$ chunk at IoU $\geq$ 0.3
  \end{itemize}
\end{column}
\begin{column}{0.48\textwidth}
  \textbf{Answer Quality}
  \begin{itemize}
    \item \textbf{LLM-Judge Score (1--5)} — GPT-4o scores generated answer
          against reference; 85\% agreement with human experts (Zheng et al.\ 2023)
  \end{itemize}

  \medskip
  \textbf{Citation Accuracy}
  \begin{itemize}
    \item Fraction of cited time spans that overlap the correct ground-truth evidence window
    \item Directly measures whether the system \emph{knows what it used}
  \end{itemize}

  \medskip
  \begin{exampleblock}{Primary metric: mean tIoU}
    At the $\pm$25.8s span annotation error, tIoU@0.3 remains a valid threshold
    (IoU $\approx$ 0.4--0.6 at that error level).
  \end{exampleblock}
\end{column}
\end{columns}

\note{The LLM-judge is GPT-4o — distinct from both the QA generator (Claude Sonnet 4.6)
and the answer generator (GPT-4o-mini). Three different models play three different roles,
which prevents circular evaluation.}
\end{frame}

%------------------------------------------------------------
\begin{frame}{Results: Single-Hop Questions ($n = 238$)}

\begin{columns}[T]
\begin{column}{0.52\textwidth}
  \small
  \begin{tabular}{lcc}
  \toprule
  Metric & Config 1 & Config 2 \\
         & (Transcript) & (+Frames) \\
  \midrule
  Mean Temporal IoU & 0.175 & \textbf{0.281} \\
  IoU@0.3           & 0.248 & \textbf{0.387} \\
  IoU@0.5           & 0.164 & \textbf{0.290} \\
  \midrule
  Hit Rate@1        & 0.218 & \textbf{0.244} \\
  Hit Rate@3        & 0.382 & \textbf{0.420} \\
  Hit Rate@5        & 0.445 & \textbf{0.529} \\
  \midrule
  LLM-Judge (1--5)  & 2.68  & \textbf{3.53} \\
  Citation Accuracy & 0.269 & \textbf{0.433} \\
  \bottomrule
  \end{tabular}
\end{column}
\begin{column}{0.44\textwidth}
  \textbf{Key observations:}
  \begin{itemize}
    \item Config 2 wins on \textbf{all 8 metrics}
    \item Largest gains: tIoU \textbf{+60\%}, citation accuracy \textbf{+61\%}
    \item LLM-judge +32\%: better retrieval $\Rightarrow$ more accurate answers
    \item Hit Rate gains are modest (+12\% at $k{=}1$): frame captions improve
          \emph{quality} of retrieved chunks, not their presence in top-$k$
  \end{itemize}

  \medskip
  135 of 238 pairs (57\%) are visual-dependent — these drive the large tIoU
  and citation accuracy gains.
\end{column}
\end{columns}

\note{[TABLE] This reproduces Table 3 (single-hop results) from Section 5.1.
Use the paper's Table 3 as a high-res screenshot for the presentation if needed.

Talking point: the modest HR@1 gain (+12\%) vs large tIoU gain (+60\%) tells an
important story — frame captions don't help the retriever FIND chunks it would have
missed, but they make the retrieved chunks CONTAIN the right evidence.}
\end{frame}

%------------------------------------------------------------
\begin{frame}{Results: Multi-Hop Questions ($n = 460$)}

\begin{columns}[T]
\begin{column}{0.52\textwidth}
  \small
  \begin{tabular}{lcc}
  \toprule
  Metric & Config 1 & Config 2 \\
         & (Transcript) & (+Frames) \\
  \midrule
  Mean Temporal IoU  & 0.180 & \textbf{0.204} \\
  IoU@0.3            & 0.274 & \textbf{0.291} \\
  IoU@0.5            & \textbf{0.076} & 0.072\textsuperscript{*} \\
  \midrule
  Hit Rate@1         & 0.326 & \textbf{0.341} \\
  Hit Rate@3         & 0.563 & \textbf{0.620} \\
  Hit Rate@5         & 0.676 & \textbf{0.722} \\
  \midrule
  LLM-Judge (1--5)   & 2.95  & \textbf{3.45} \\
  Citation Accuracy  & 0.273 & \textbf{0.333} \\
  \bottomrule
  \end{tabular}

  \footnotesize\textsuperscript{*}Frame-aug chunks are longer, widening the union
  span and reducing tIoU even when the correct chunk is retrieved.
\end{column}
\begin{column}{0.44\textwidth}
  \textbf{Key observations:}
  \begin{itemize}
    \item Config 2 wins 7/8 metrics
    \item \textbf{Smaller gains} than single-hop (+13\% tIoU vs +60\%)
    \item Mixed composition explains this: 320/460 are visual pairs (benefit from
          captions), 140/460 are text-only (slight regression)
    \item Hit Rate@5 is 107\% higher than Hit Rate@1 — multi-hop retrieval
          is \textbf{$k$-sensitive}: larger $k$ budgets would help
    \item One IoU@0.5 exception: larger union span from longer chunks
  \end{itemize}
\end{column}
\end{columns}

\note{[TABLE] This reproduces Table 4 (multi-hop results) from Section 5.2.

The IoU@0.5 footnote is important: explain that when more chunks are cited
(because captions make more of them seem relevant), the union span gets wider,
which hurts precision but doesn't hurt Hit Rate since that only checks presence.}
\end{frame}

%------------------------------------------------------------
\begin{frame}{Key Finding: Visual vs.\ Text-Only Split}

\begin{center}
\small
\begin{tabular}{lcccccc}
\toprule
 & \multicolumn{3}{c}{\textbf{Config 1: Transcript-Only}}
 & \multicolumn{3}{c}{\textbf{Config 2: Transcript+Frames}} \\
\cmidrule(lr){2-4}\cmidrule(lr){5-7}
Question type & tIoU & Judge & Cit.Acc & tIoU & Judge & Cit.Acc \\
\midrule
Visual-dep.\ ($n=455$) & 0.129 & 2.57 & 0.198
                        & \cellcolor{green!20}\textbf{0.213} & \cellcolor{green!20}\textbf{3.52} & \cellcolor{green!20}\textbf{0.340} \\
Text-only ($n=243$)    & \textbf{0.271} & 3.38 & 0.409
                        & 0.263 & \textbf{3.41} & \textbf{0.418} \\
\bottomrule
\end{tabular}
\end{center}

\bigskip
\begin{columns}[T]
\begin{column}{0.48\textwidth}
  \textbf{Visual-dependent questions:}\\
  Frame captions yield \textbf{+65\% tIoU} (0.129 $\to$ 0.213),
  \textbf{+37\% LLM-judge}, and \textbf{+72\% citation accuracy}
\end{column}
\begin{column}{0.48\textwidth}
  \textbf{Text-only questions:}\\
  Essentially \textbf{unchanged} ($-$3\% tIoU, $+$1\% judge):
  frame captions add \textbf{no noise} and \textbf{no benefit}
  where visual evidence is not needed
\end{column}
\end{columns}

\medskip
\begin{center}
\textbf{\large Frame captions add signal only where visual evidence is required.}
\end{center}

\note{[TABLE] This is Table 5 from Section 5.2 (visual\_vs\_text). This is the most
important table in the paper — use it directly as a screenshot for maximum impact.

This is the central empirical claim of the paper. The isolation is clean:
+65\% on the group that needs visual evidence, essentially 0\% change on the
group that doesn't. This demonstrates both (a) that frame captions work and
(b) that they don't introduce noise into text-based retrieval.}
\end{frame}

%------------------------------------------------------------
\begin{frame}{Failure Mode Analysis}

\begin{columns}[T]
\begin{column}{0.45\textwidth}
  \textbf{Three recurring failure modes:}

  \medskip
  \textbf{1. Span inflation}\\
  Entire 45s chunk cited when relevant evidence occupies only a small portion
  $\Rightarrow$ tIoU penalty even on correct retrieval

  \medskip
  \textbf{2. Visual blindness (Config 1)}\\
  Answer appears only on a lecture slide — absent from spoken transcript.
  Config 2 recovers by reading frame caption.\\
  {\small Mean tIoU: 0.129 $\to$ 0.213 (+65\%)}

  \medskip
  \textbf{3. Partial multi-hop retrieval}\\
  Retriever surfaces one hop at rank 1 but misses the second.\\
  Hit Rate@5 is \textbf{107\% higher} than Hit Rate@1 on multi-hop.

\end{column}
\begin{column}{0.52\textwidth}
  \small
  \setlength{\fboxsep}{6pt}
  \fbox{
  \begin{minipage}{\textwidth}
  \textbf{Visual-blindness case} (\texttt{nyu\_dl\_week10\_q008})

  \smallskip
  \textbf{Q:} What are the three bullet points on the `Why self-supervision?' slide?

  \smallskip
  \textbf{Config 1} top chunk: spoken transcript only.\\
  \textit{Answer:} ``The provided excerpts do not contain enough information.''\\
  LLM-judge: \textbf{1.5/5}

  \smallskip
  \textbf{Config 2} same chunk, with frame caption:\\
  \texttt{[frame: Slide `Why self-supervision?' --- (1) Helps us learn using
  observations; (2) Does not require exhaustive annotation; (3) Leverage multiple
  modalities]}\\
  \textit{Answer:} lists all three bullets correctly.\\
  LLM-judge: \textbf{5.0/5}
  \end{minipage}
  }

  \smallskip
  \footnotesize tIoU: 0.000 (Config 1) $\to$ 1.000 (Config 2)
\end{column}
\end{columns}

\note{[FIGURE] This reproduces the content of Figure 5 (failure case, fig:failure\_case)
from Section 5.3. The boxed example is already included on this slide, but you can
swap it for the actual Figure 5 screenshot from the paper for consistency.

The visual-blindness case is the paper's most concrete and compelling example.
Both configs retrieve the SAME top chunk — the only difference is that Config 2's
chunk contains the slide text. This eliminates any confound from retrieval differences.}
\end{frame}

%------------------------------------------------------------
\begin{frame}{Hit Rate Analysis: $k$-Sensitivity of Multi-hop Retrieval}

\begin{columns}[T]
\begin{column}{0.44\textwidth}
  \begin{center}
  \fbox{\rule{0pt}{5.0cm}\rule{4.8cm}{0pt}}\\
  \footnotesize\textit{[Insert Figure 4 — fig\_hr\_at\_k]}
  \end{center}
\end{column}
\begin{column}{0.52\textwidth}
  \textbf{Reading the plot:}
  \begin{itemize}
    \item Solid lines = single-hop; dashed = multi-hop
    \item Config 2 (+Frames) consistently above Config 1 across all $k$
  \end{itemize}

  \medskip
  \textbf{Key insight — multi-hop curves:}
  \begin{itemize}
    \item Steep rise from $k{=}1$ to $k{=}5$: the second hop is often ranked 2--5
    \item Config 1 multi-hop: HR@1 = 0.326, HR@5 = 0.676 (\textbf{+107\%})
    \item Config 2 multi-hop: HR@1 = 0.341, HR@5 = 0.722 (\textbf{+112\%})
    \item $\Rightarrow$ increasing retrieval budget $k$ would directly improve multi-hop performance
  \end{itemize}

  \medskip
  \textbf{Single-hop curves} plateau earlier — first hop dominates at low $k$.
\end{column}
\end{columns}

\note{[FIGURE] Place Figure 4 (images/fig\_hr\_at\_k.pdf) in the left placeholder.
This is the HR@k line plot from Section 5.2 (fig:hr\_at\_k).

Key talking point: the steep multi-hop curve tells us that the retriever IS finding
the second hop — it just lands at rank 3--5, not rank 1. A re-ranking step or
larger $k$ budget would likely close most of the remaining gap on multi-hop.}
\end{frame}

%------------------------------------------------------------
\begin{frame}{Summary of Results}

\begin{center}
\small
\begin{tabular}{lcccc}
\toprule
Metric & Config 1 & Config 2 & $\Delta$ & Where the gain comes from \\
\midrule
Mean tIoU (all)         & 0.154 & \textbf{0.198} & +29\% & Visual-dep.\ +65\%, text 0\% \\
IoU@0.3                 & 0.228 & \textbf{0.279} & +22\% & \\
Hit Rate@5              & 0.515 & \textbf{0.565} & +10\% & \\
LLM-Judge (1--5)        & 3.03  & \textbf{3.56}  & +17\% & \\
Citation Accuracy       & 0.234 & \textbf{0.316} & +35\% & \\
\midrule
tIoU — visual-dep.\ ($n$=455)  & 0.129 & \textbf{0.213} & \textcolor{highlight}{\textbf{+65\%}} & Frame caps expose slide text \\
tIoU — text-only ($n$=243)     & \textbf{0.271} & 0.263 & $-$3\% & No visual evidence needed \\
\bottomrule
\end{tabular}
\end{center}

\medskip
\begin{columns}[T]
\begin{column}{0.48\textwidth}
  \begin{exampleblock}{What works}
    Frame captions close the visual-evidence gap precisely — +65\% tIoU and +72\%
    citation accuracy on visual-dependent questions.
  \end{exampleblock}
\end{column}
\begin{column}{0.48\textwidth}
  \begin{alertblock}{What does not help}
    Frame captions do not degrade transcript-only retrieval ($-$3\% tIoU,
    within noise). No negative transfer observed at the aggregate level.
  \end{alertblock}
\end{column}
\end{columns}

\note{This is a synthesis slide with no direct paper figure — build from the numbers
in Tables 3, 4, and 5. The $\Delta$ column is the paper's abstract numbers.

Use this as your closing results slide before the conclusion.}
\end{frame}

%------------------------------------------------------------
\begin{frame}{Limitations}

\begin{itemize}
  \setlength{\itemsep}{8pt}
  \item \textbf{English-only:} corpus limited to English lectures with YouTube
        auto-captions; low-quality VTT files introduce occasional transcription
        errors in technical terms.

  \item \textbf{LLM-only QA review:} ground-truth spans are estimated from chunk
        boundaries ($\pm$15--30s) rather than human-tightened to exact evidence windows.
        Mean annotation error $\approx$25.8s (worst case 82.5s).

  \item \textbf{Frame captioner not fine-tuned:} Qwen2-VL-7B occasionally misreads
        small mathematical notation on lecture slides.

  \item \textbf{Scale:} 810 pairs across 60 lectures is smaller than existing benchmarks
        (EduVidQA: 5,252; LongVidSearch: 3,000), justified by higher per-pair curation
        cost (temporal grounding, multi-hop validation, cross-family review).

  \item \textbf{Single embedding model:} results are specific to all-MiniLM-L6-v2;
        stronger embeddings may compress the gap.
\end{itemize}

\note{Be prepared for the advisor to ask about the scale limitation. The response is:
(1) 12 pairs/lecture exceeds LongVidSearch (6.7/lecture) in density; (2) Min et al.
(2019) show that strict non-adjacency enforcement is more important than raw pair count
for genuine multi-hop validity; (3) 3,000 pairs is feasible in ~$55 and ~5 hours
of API calls if needed.}
\end{frame}

%------------------------------------------------------------
\begin{frame}{Future Work}

\begin{enumerate}
  \setlength{\itemsep}{10pt}
  \item \textbf{Non-English lectures:} extend the corpus to multi-lingual content
        and evaluate ASR quality as a variable.

  \item \textbf{Semantic chunking:} replace sliding-window with topic-boundary detection
        to produce semantically coherent chunks.

  \item \textbf{Joint text-visual embedding:} train a lecture-specific encoder that
        directly maps frame+caption pairs into a shared retrieval space with transcript text.

  \item \textbf{Stronger generators:} evaluate GPT-4o and Claude Opus on the benchmark
        to measure the ceiling of RAG-constrained generation.

  \item \textbf{Selective frame augmentation:} attach captions only to chunks where
        visual evidence is detected — avoids the $-$5.5\% tIoU signal dilution observed
        on transcript-only multi-hop pairs.
\end{enumerate}

\note{Item 5 (selective augmentation) is the most actionable near-term improvement —
it directly addresses the one case where Config 2 underperforms Config 1 (text-only
multi-hop pairs). The basic approach would be a visual-activity detector run on frames
before the captioning step.}
\end{frame}

%------------------------------------------------------------
\begin{frame}{Conclusion}

\begin{columns}[T]
\begin{column}{0.55\textwidth}
  \textbf{What we built:}
  \begin{itemize}
    \item LectureBench — 810 temporally grounded QA pairs over 60 long lectures
    \item Only benchmark combining multi-hop + temporal spans + visual split for lecture-length video
    \item Open-source RAG pipeline with two configurations and full evaluation suite
  \end{itemize}

  \medskip
  \textbf{What we found:}
  \begin{itemize}
    \item Frame captions: \textbf{+65\% tIoU} on visual questions
    \item Text-only questions: \textbf{$-$3\%} (no degradation)
    \item Multimodal augmentation is \textbf{surgically precise}
    \item Multi-hop retrieval is \textbf{$k$-sensitive} — larger budgets would help
    \item Visual blindness is a real and measurable failure mode in transcript-only RAG
  \end{itemize}
\end{column}
\begin{column}{0.42\textwidth}
  \begin{exampleblock}{The one-slide answer}
    Adding Qwen2-VL-7B frame captions to lecture chunks improves temporally
    grounded retrieval by \textbf{+65\%} on visual-dependent questions
    and leaves text-only questions unchanged, directly confirming that
    transcript-only RAG has a structural visual-evidence blind spot.
  \end{exampleblock}

  \bigskip
  \textbf{Code \& data:}\\
  \small\texttt{github.com/NassimF/VQA-Benchmark}\\
  CC BY-NC-SA 4.0
\end{column}
\end{columns}

\note{Final talking point: the paper's value is not just the benchmark numbers —
it is the controlled experimental design (single variable: chunk content) that
makes the causal claim about visual evidence defensible. No joint embedding,
no fine-tuning, no architecture change — just what you put in the text field.}
\end{frame}

%------------------------------------------------------------
\begin{frame}[plain]
\begin{center}
  {\Large\bfseries Thank You}\\[12pt]
  {\large Questions?}\\[24pt]
  \texttt{seyedehnasim.faridnia@utsa.edu}\\[6pt]
  \texttt{github.com/NassimF/VQA-Benchmark}
\end{center}
\end{frame}

%=============================================================
\end{document}
%=============================================================
